%!TEX root = ../thesis.tex

本研究では，深層強化学習を用いたヒューマノイドロボットの歩行動作を対象とした，
報酬関数の設計が歩行性能や動作に与える影響を明らかにするため，歩行速度追従率，胴体の姿勢偏差，床反力，
歩幅および左右対称性，重心位置と足上げ高さ，各関節トルクと角速度の時系列グラフを用いて複数の
定量的な評価を行った．\\
目的に応じて報酬項のバランスを調整にする必要がある

\section{joint deviation(hip)の影響}
本節は，\textbf{joint deviation(hip)}について議論する．

\section{Lin Vel Z L2の影響}
本節は，\textbf{Lin Vel Z L2}について議論する．

\section{DOF Acc L2の影響}
本節は，\textbf{DOF Acc L2}について議論する．

\section{DOF Torques L2の影響}
本節は，\textbf{DOF Torques L2}について議論する．

\section{Feet Slideの影響}
本節は，\textbf{Feet Slide}について議論する．

\section{Flat orientationの影響}
本節は，\textbf{Flat orientation}について議論する．